\documentclass[11pt,a4paper,landscape]{article}
\usepackage[T1]{fontenc}\usepackage[utf8]{inputenc}\usepackage[french]{babel}
\usepackage{lmodern,amsmath,amssymb,graphicx,booktabs,xcolor,hyperref,fancyhdr}
\usepackage[margin=16mm]{geometry}
\hypersetup{colorlinks=true,linkcolor=blue,urlcolor=blue,pdftitle={Plaque trouée : champ dense et remapping des images}}
\graphicspath{{../results/hudge_dense_remap_20261115/figures/}}
\pagestyle{fancy}\fancyhf{}\fancyhead[L]{WGPU-Matching — Initialisation et images remappées}
\fancyhead[R]{15 novembre 2026}\fancyfoot[R]{\thepage}\setlength{\headheight}{14pt}
\setlength{\parindent}{0pt}\setlength{\parskip}{5pt}
\newcommand{\fig}[2]{\begin{center}\includegraphics[width=\linewidth,height=#2,keepaspectratio]{#1}\end{center}}
\begin{document}
{\LARGE\bfseries Plaque trouée : reconstruire une initialisation\par
et voir réellement les images remappées\par}
\vspace{3mm}
\textbf{Rectification après examen des textures : la recommandation comme initialisation près des trous est retirée.}
La correction des contours rapproche les grands trous de leur forme de référence, mais dégrade
le mouchetis entre 4 et 16 pixels des trous. Dans les bandes 4--8 puis 8--16 pixels, la corrélation
des intensités passe respectivement de 0,700 à 0,614 et de 0,864 à 0,774. L'accord des contours
ne valide pas la correspondance matérielle. Les images de ce document restent des diagnostics.

Sur 825 centres décalés de ceux ayant servi à construire le champ, le solveur local retient
764 fenêtres avec le champ grossier, 780 avec le champ dense, puis 782 avec les contours.
\textbf{Aucun de ces centres n'est à moins de 8 pixels d'un trou.} Ces comptes ne qualifient donc pas
la couronne proche des trous. Il n'y a pas de vérité de déplacement associée à ces images.

\fig{native_images.pdf}{100mm}
\textbf{Les trois images ont exactement le même canevas de 1024\,$\times$\,1024 pixels.}
La matière de référence n'en occupe qu'une partie. L'image corrigée est obtenue par échantillonnage
bilinéaire de la déformée originale. Aucun trou n'est redessiné, aucun cercle n'est imposé,
aucun masque de référence n'est utilisé pour repeindre les pixels.

\newpage
\section*{1. Les images, aux mêmes coordonnées de référence}
\fig{remapping_stages.pdf}{147mm}
Chaque ligne compare une paire indépendante à la même référence, sans propagation temporelle.
Les recadrages sont identiques ; ils ne modifient ni les coordonnées ni le champ exporté.
Les ancrages sont les sorties retenues des fenêtres de 16 pixels, espacées de 8 pixels.

\newpage
\section*{2. Ce qui est corrigé, et ce qui ne l'est pas}
\fig{holes_191.pdf}{120mm}
Le rapport d'axes est obtenu à partir des moments d'aire du trou segmenté : il vaut un pour un
cercle idéal, mais ne mesure pas à lui seul l'accord de taille ou de position.
Au niveau 191, le grand trou central atteint un rapport de 1,014, un recouvrement de 97,8\,\%
et un écart de centre de 0,10 pixel. Le petit trou proche de $(265,101)$ reste à un rapport
de 1,479 et seulement 66,3\,\% de recouvrement. La forme et la position doivent être examinées ensemble.

\textbf{Limite essentielle :} les contours participent à l'ajustement ; leur recouvrement n'est
donc pas une mesure indépendante de précision mécanique. La corrélation globale du mouchetis
passe de 0,806 à 0,828 au niveau 191. Une silhouette mieux recalée ne suffit pas à prouver
un déplacement plus exact dans la matière.

\newpage
\section*{3. Construction du champ et signification des gradients}
La corrélation d'images numériques (DIC) est raffinée ici par le solveur
Gauss--Newton inverse compositionnel (IC-GN). Chaque fenêtre résolue fournit un déplacement
$\mathbf d_i$ et son gradient $G_i$, exprimés en pixels natifs.
Le champ dense initial combine les développements locaux :
\[
 \mathbf u_0(\mathbf x)=\sum_i\alpha_i(\mathbf x)
 [\mathbf d_i+G_i(\mathbf x-\mathbf x_i)],\qquad
 \alpha_i=\frac{\exp(-\|\mathbf x-\mathbf x_i\|^2/(2\ell^2))}
 {\sum_j\exp(-\|\mathbf x-\mathbf x_j\|^2/(2\ell^2))},\quad \ell=12\ \mathrm{px}.
\]
La somme utilise les 24 voisins les plus proches. Les dérivées des poids sont incluses
dans $\nabla\mathbf u_0$ ; les quatre gradients ne sont pas interpolés indépendamment du déplacement.
La troncature aux voisins rend ce champ seulement lisse par morceaux.

Les contours sont extraits des grandes composantes du fond uniforme réellement observé
dans ces rendus. Une correction par fonctions de base radiale (RBF) gaussiennes est ajoutée :
\[
 \boldsymbol\varphi(\mathbf x)=\mathbf x+\mathbf u_0(\mathbf x)+
 \sum_{k=1}^{384}\mathbf c_k K_k(\mathbf x),\qquad
 K_k=\exp[-\|\mathbf x-\mathbf s_k\|^2/(2\cdot12^2)].
\]
Avec $D_1$ la distance signée au contour de la cible et $\Gamma_0$ le contour de référence,
les coefficients minimisent, par six linéarisations successives,
\[
 4\sum_{\mathbf x\in\Gamma_0}D_1(\boldsymbol\varphi(\mathbf x))^2+
 \sum_{i\in\mathrm{ancrages}}\|\boldsymbol\varphi(\mathbf x_i)-
 (\mathbf x_i+\mathbf u_0(\mathbf x_i))\|^2+0{,}05\sum_k\|\mathbf c_k\|^2.
\]
La linéarisation du premier terme est $D_1+\nabla D_1^{\mathsf T}\delta\boldsymbol\varphi$ :
elle contraint la normale au contour, pas une correspondance tangentielle arbitraire.
Une recherche de pas refuse les corrections créant un déterminant de jacobienne non positif
sur les pixels de matière. Elle ne borne pas l'amplitude du gradient. Au niveau 150,
plusieurs corrections sont refusées : le contrôle n'est pas contourné pour améliorer la figure.

\newpage
\section*{4. Le champ est-il une meilleure initialisation ?}
\fig{offset_validation.pdf}{92mm}
Les centres de ce contrôle sont décalés de $(4,4)$ pixels par rapport à la grille des ancrages.
Les fenêtres font 24\,$\times$\,24 pixels, avec un pas de 8. Leurs supports recouvrent les observations
utilisées dans le champ : ce n'est pas une validation sur de nouvelles images.

\begin{center}\begin{tabular}{lrrr}\toprule
Initialisation & Statut convergé après 30 itérations & Retenues après continuation & Total\\\midrule
Champ grossier, 36 ancrages &699&764&825\\
Champ dense de fenêtres 16 pixels &734&780&825\\
Champ dense et contours, rayon 16, 196 bases &727&784&825\\
Champ dense et contours, rayon 12, 384 bases &729&782&825\\\bottomrule
\end{tabular}\end{center}
La continuation comporte au plus quatre appels de 30 itérations ; seules les sorties non convergées
et finies sont relancées. Le seuil d'arrêt vaut $10^{-3}$ et la corrélation normalisée minimale 0,65.
Une fenêtre retenue doit en outre rester dans l'image et conserver un déterminant positif sur son support.
Les statuts seuls donnent respectivement 767, 783, 786 et 784 candidats après continuation ; les contrôles
finaux expliquent les effectifs retenus plus petits. Les temps archivés ne constituent pas un benchmark :
ils incluent notamment des coûts de première compilation différents.

Le champ est exporté aux centres sous l'ordre
$(u,\partial_xu,\partial_yu,v,\partial_xv,\partial_yv)$, avec les pixels natifs comme unités.
Les fichiers distinguent les valeurs prédites par le champ et les paramètres locaux effectivement mesurés.

\newpage
\section*{5. Bilan pratique et expériences conservées}
\textbf{Choix pour cet essai.} Utiliser les 36 ancrages grossiers pour initialiser une grille plus fine,
résoudre les fenêtres de 16 pixels, reconstruire le champ, puis exploiter les contours observables
pour améliorer les zones proches des trous. Au niveau 191, une réinjection du champ corrigé
augmente les ancrages fins de 784 à 790 sur 979. Les 189 fenêtres restantes ne sont pas déclarées résolues.

\textbf{Résultats négatifs.} Les champs obtenus à partir des fenêtres quadratiques de 32 pixels
recalent moins bien l'image entière dans ce protocole, malgré leur intérêt photométrique local.
Les fenêtres quadratiques de 16 pixels et les fenêtres affines de 8 pixels produisent des champs
avec replis lors de leur combinaison ; ces champs ne sont pas retenus pour l'export recommandé.
Une RBF globale ajustée aux ancrages denses lisse trop les variations. Elle ne remplace pas
le traitement des contours. Une correction de flot photométrique exploratoire améliore la corrélation,
mais ses répétitions créent des replis ; elle n'entre pas dans le résultat livré.

\textbf{Limites scientifiques.} Ces résultats concernent trois niveaux d'une même plaque.
La segmentation par fond uniforme est spécifique à ces rendus ; elle ne constitue pas un segmentateur
général de fissures ni un protocole validé sur des images expérimentales. Les déplacements du maillage
en croix fourni dans l'autre dossier ne correspondent pas à cette plaque et ne sont pas utilisés.
Sans vérité de déplacement associée, ni le statut du solveur ni le bon accord des trous ne certifient
une erreur mécanique donnée. La petite ouverture reste une zone d'échec explicite.

\textbf{Exécution.} Les résolutions locales utilisent WebGPU, interface portable de calcul sur
processeur graphique. L'ajustement du champ, les distances aux contours et le remapping de ce
prototype sont calculés sur le processeur central. Cette expérience ne prétend donc ni être entièrement
résidente sur processeur graphique, ni tenir une cadence temps réel. Aucun code du solveur externe
ni aucune conclusion du manuscrit principal n'ont été modifiés.

\textbf{Fichiers pour inspection.} Pour chaque niveau, le dossier
\texttt{results/hudge\_dense\_remap\_20261115/refined/level\_XXX/} contient
la référence, la déformée originale, la déformée corrigée, une animation alternée,
les déplacements et gradients natifs, les initialisations aux centres et les statistiques par trou.
La galerie \texttt{index.html} à la racine du résultat donne accès aux images natives et aux animations.

\textbf{Conclusion.} Les remappings rendent enfin visible l'effet du champ non affine.
La densification améliore l'initialisation ; les contours améliorent surtout la géométrie visible.
L'image corrigée est utile pour analyser les échecs, mais ne doit pas être présentée comme une
mesure DIC achevée ni comme une restitution parfaite de tous les trous.
\end{document}
